\section{Level structures, sheaves, and specialization}
\label{found:section}

The construction will pass through enhanced flags and through a quotient
that identifies paired enhancements on two marked curves.  We describe these stacks
first, then fix the sheaf and Hecke conventions that allow the resulting
objects to be specialized.  The nearby-cycle inputs are
\cite[Propositions~2.1 and~5.1]{CT}; their hypotheses concern the local
geometry at the moving legs, rather than the number of level points.

\subsection{The level stacks}
\label{found:level-stacks}

Let $X$ be a smooth projective connected curve over an algebraically
closed field, and let $D=x_1+\cdots+x_s$ be a reduced divisor.  Write
$U=X\setminus |D|$.  We use $G=\mathrm{SL}_n$, the upper triangular
Borel $B$, its unipotent radical $N_B$, and the diagonal torus
$T=B/N_B$.  The stacks
\[
 \cA=\Bun_{G,B,D}(X),\qquad
 \cA^+=\Bun_{G,N_B,D}(X)
\]
classify a $G$-bundle with, respectively, a $B$-reduction or an
$N_B$-reduction at every $x_i$.  We call the latter an enhanced flag.
Forgetting the enhancements gives a representable torsor
\begin{equation}\label{found:torsor}
                         q:\cA^+\longrightarrow\cA
                         \quad\text{under }T^s.
\end{equation}
Unlevelled bundles and opposite Borels are also allowed.  All these
stacks are smooth and locally of finite type: the bundle deformation
obstruction lies in $H^2$ on a curve, and adding a reduction at a point
has smooth fiber $G/B$ or $G/N_B$.

For two such pointed curves, identify the flag tori at corresponding
points by fixed isomorphisms $\tau_i:T_{1i}\simeq T_{2i}$.  Put
\begin{equation}\label{found:quotient}
 \cB_{12}=[(\cA_1^+\times\cA_2^+)/T^s_\Delta],\qquad
 T^s_\Delta=\prod_{i=1}^s\{(t,\tau_i(t)):t\in T_{1i}\}.
\end{equation}
The branch convention in the nodal construction will specify the
$\tau_i$.  This smooth stack maps to $\cA_1\times\cA_2$ as a torsor
under $(T_1^s\times T_2^s)/T^s_\Delta$.  We will use this quotient in
characteristic zero.  Thus its objects are pairs of enhanced bundles
modulo simultaneous changes of the two enhancements at each paired
marking.

A Hecke modification in $U$ transports every level structure, because
it identifies the bundles near all the $x_i$.  On the product and the
quotient, a leg in $U_1\sqcup U_2$ modifies the corresponding component.
The quotient Hecke correspondence pulls back to the product
correspondence when either its input or its output is lifted to the
product.  Indeed the transported enhancements provide the lift on the
other side.  This cartesian property also holds for successive
modifications and their collision diagrams.

On a fixed Schubert bound, both the input-and-leg and output-and-leg
maps are representable and proper.  On an exact-position stratum they
are smooth.  To see that adding all the levels preserves these facts,
choose a coordinate and a frame on the formal disc at the leg.  The
modification is a split affine-Grassmannian model there; the marked
fibers remain in the complementary piece of the gluing.  For a fixed
bound the frame can be taken to finite jet order, so these descriptions
are obtained by smooth covers.  The same description applies to a leg
in the smooth scheme locus of a family of twisted nodal curves.  Its
stacky nodes remain in the complementary piece.  Restricting to fixed
node types restricts both bounded projections by base change and hence
preserves their properness.  On a finite-type output chart a fixed
bound meets only a quasi-compact part of the input stack.

\subsection{Local constructibility and relative Satake}

Put $E=\overline{\mathbb Q}_\ell$.  For a smooth stack $\cB$ locally of
finite type, $\Dlc(\cB,E)$ denotes the geometric adic complexes whose
ordinary pullback to every smooth finite-type scheme chart is bounded
constructible.  Thus ``lcc'' means locally bounded constructible, and
bounds may depend on the chart.  An object $F$ is perverse when
$f^*F[d]$ is perverse for every smooth chart $f:S\to\cB$ of constant
relative dimension $d$.  These conditions define a perverse
$t$-structure and its truncations locally on the stack.

For a closed conic subset $C\subset T^*\cB$, set
\[
 f^\circ C=\operatorname{image}
   (S\times_{\cB}C\longrightarrow T^*S).
\]
The assertion $\SS(F)\subset C$ means
$\SS(f^*F)\subset f^\circ C$ on every such chart.  Here stack
cotangent vectors mean degree-zero cotangent vectors, and singular
support is the geometric adic singular support, with the
function-test characterization of \cite[\S1.5, Theorem~(vii)]{Barrett}.
Shifts do not affect singular support.

Fix a Satake equivalence with $\Rep_E(\Gd)$, where $\Gd=\mathrm{PGL}_n$,
including its fusion commutativity constraint.  Write $\rho_G$ for the
half-sum of positive roots.  For a dominant
coweight $\lambda$ its simple kernel has open-orbit restriction
\begin{equation}\label{found:satake-normalization}
 \IC_\lambda|_{\Gr^\lambda}
       =E[d_\lambda](d_\lambda/2),\qquad
 d_\lambda=\langle 2\rho_G,\lambda\rangle.
\end{equation}
The integer $d_\lambda$ is even: the coweight lattice of $\mathrm{SL}_n$
is the coroot lattice, and $\langle2\rho_G,\alpha^\vee\rangle=2$ for a
simple coroot.  Tate twists here are geometric coefficient lines;
compatible trivializations may be chosen for Betti comparison.

For a finite set $I$ let $\cH_I$ be the moving Hecke stack, with input
map $p_I$ and output-and-leg map $o_I:\cH_I\to\cB\times U^I$.
The relative Satake kernel $\operatorname{Sat}_I(\boldsymbol V)$ for
$\boldsymbol V=(V_i)_{i\in I}$ is the exterior product of the
fixed-point kernels at distinct legs, extended by fusion at collisions.
Equivalently, it is computed using successive modifications and proper
convolution.  We define
\begin{equation}\label{found:hecke}
 \Hecke_{I,\boldsymbol V}(F)
   =o_{I,*}\bigl(F\,\widetilde\boxtimes\,
                         \operatorname{Sat}_I(\boldsymbol V)\bigr).
\end{equation}
In input frames the twisted product is the ordinary exterior product;
Satake equivariance supplies descent.  The push is taken on a finite
bound containing the kernel support.  There is no moving-leg shift in
\eqref{found:hecke}; in particular
$\Hecke_{I,(\mathbf1)}(F)=F\boxtimes E_{U^I}$.

\begin{definition}\label{found:coherence}
A coherent eigenstructure for a tensor local system
$V\mapsto L_V$ on $U$ consists of natural isomorphisms
\[
 \Hecke_{I,\boldsymbol V}(F)
       \simeq F\boxtimes\Bigl(\mathop{\boxtimes}_{i\in I}L_{V_i}\Bigr)
\]
for all finite sets and representations, compatible with the unit,
successive modifications and convolution, permutations, and fusion.
Explicitly, if $a:I\twoheadrightarrow J$ and
$\delta_a:U^J\to U^I$ is its collision diagonal, the fusion comparison
uses ordinary pullback and identifies
\[
 (1\times\delta_a)^*\Hecke_{I,\boldsymbol V}(F)
 \simeq\Hecke_{J,\boldsymbol W}(F),\qquad
 W_j=\bigotimes_{i\in a^{-1}(j)}V_i.
\]
The eigenmaps must commute with these comparisons, including their
compositions for nested collisions.
\end{definition}

We use this convention for products and quotients as well, allowing
legs on both component curves.  When studying perversity over $U^I$
we will insert $[|I|]$ explicitly and subsequently remove it.

\subsection{Geometric nearby cycles}

Let $R$ be an excellent complete strictly henselian discrete valuation
ring with algebraically closed residue field and with $\ell$ invertible.
Write $S=\Spec R$, fix a geometric generic point $\bar\eta$, and write
$0$ for the closed point.  Geometric nearby cycles
$\Psi_Y:D^b_c(Y_{\bar\eta},E)\to D^b_c(Y_0,E)$ take no inertia
invariants.  If descent data supply an inertia action, we use the
underlying geometric complex.

\begin{proposition}[Nearby-cycle input]\label{found:nearby}
For finite-type $S$-schemes, geometric nearby cycles preserve bounded
constructibility, are perverse exact, commute with smooth pullback and
proper pushforward, and satisfy exterior-product comparison.  They are
unchanged by finite extensions of the trait inside the fixed geometric
generic field.  The same statements hold locally on the smooth stacks
used here, with representable proper maps.

Let $Z$ be a finite-type $S$-scheme.  Suppose a lisse finite-rank sheaf
$L_{\bar\eta}$ on $Z_{\bar\eta}$ has
a lattice over the integers of a finite extension of $\mathbb Q_\ell$.
Assume that each finite reduction extends lisse over $Z$ after a finite
extension of $S$, compatibly on further common extensions, with special
reductions belonging to a fixed lisse sheaf $L_0$ on $Z_0$.  The trait
extension may depend on the reduction.  Then for any map $g:Y\to Z$
of finite-type $S$-schemes and any bounded constructible $F$ on $Y_{\bar\eta}$,
\begin{equation}\label{found:lattice}
 \Psi_YF\otimes g_0^*L_0
      \xrightarrow{\sim}
       \Psi_Y(F\otimes g_{\bar\eta}^*L_{\bar\eta}).
\end{equation}
This is natural in maps and tensor products of the lisse systems, and
holds locally on the smooth stacks in use.
\end{proposition}

These are precisely the assertions of \cite[Proposition~2.1]{CT}.
We recall why they apply to geometric coefficients without a common
finite field of descent.  At a fixed torsion coefficient level, each
finite diagram descends after a finite trait extension.  The comparisons
are compatible with further extensions and with reduction of the
coefficients.  Passing to the adic system, then inverting the coefficient
uniformizer, gives the geometric functor.  Equivalently, over the
integral closure of $R$ in its geometric generic field, the universally
locally acyclic extension formalism gives nearby cycles by special
restriction; see \cite[Theorems~4.1 and~6.7]{HS}.  Smooth descent gives
the stack assertions, including perversity with the chart shift $[d]$.
For \eqref{found:lattice}, apply the lisse tensor formula to each
extended finite reduction and pass to the limit.  This permits arbitrary
$g$: it is the lisseness of the factor, rather than a base-change claim
for $\Psi$ along $g$, that is used.  No single finite extension is
required for the entire lattice tower.

All these comparisons use the natural pull, tensor, and proper-push
exchange maps.  They respect composition, tensor associativity,
projection formula, and proper base change.  We will also use the
constant-family property: a complex pulled back from a fixed geometric
fiber has zero vanishing cycles in a smooth constant family.

\begin{proposition}[Specialization of eigenstructures]
\label{found:hecke-specialization}
Let $\cB/S$ be smooth and locally of finite type, and let
$\mathscr U/S$ be a smooth curve of allowed legs.  Assume the split
spherical disc models, proper bounds, and smooth exact-position maps
described in Section~\ref{found:level-stacks}.  Then there are natural
comparisons
\[
 \Hecke^0_{I,\boldsymbol V}(\Psi F)
   \xrightarrow{\sim}
 \Psi_{\cB\times_S\mathscr U^I}
             \Hecke^{\bar\eta}_{I,\boldsymbol V}(F)
 \qquad(F\in\Dlc(\cB_{\bar\eta},E)),
\]
compatible with the entire relative Hecke system of
Definition~\ref{found:coherence}.  Consequently, if $F$ has a coherent
lisse eigenvalue whose lattices satisfy Proposition~\ref{found:nearby}
compatibly as a tensor system, $\Psi F$ has the specialized coherent
eigenvalue.
\end{proposition}

\begin{proof}
This is \cite[Proposition~5.1]{CT}; we explain the local hypothesis and
the collision assertion needed here.  In a coordinate-and-input-frame
cover a bounded step is a product with a split Schubert model.  For an
arbitrary preceding space $Y$ and a complex $K$ on it, the comparison
for the next twisted product is therefore the exterior-product map
\[
 \Psi_YK\boxtimes\Psi_{\Gr}\operatorname{Sat}(V)
 \xrightarrow{\sim}
 \Psi_{Y\times_S\Gr}(K\boxtimes\operatorname{Sat}(V)).
\]
The first factor need not be smooth.  The specialized Grassmannian
kernel is the Satake kernel with the same label and normalization;
\cite[Lemma~5.3]{CT} fixes a single tensor identification for these
kernels.  Frame descent and proper push give the one-step comparison.

For successive steps, take $Y$ to retain the preceding modifications
and all leg parameters.  The same calculation is valid after setting
any collection of legs equal.  Proper convolution then gives the fused
comparison, and naturality of the exchange maps identifies it with the
comparison obtained by successive pushes.  Thus the diagonal exchange,
which need not be invertible on arbitrary sheaves, is invertible on
these Hecke constructions and respects nested diagonals.  Permutations
are the fusion symmetry of the kernels, and the identity-modification
kernel gives the unit.  Transported levels at all marked points, and
node types away from the disc, change none of these frame calculations.
This verifies the application of the cited proposition to every stack
and family used here.  Finally \eqref{found:lattice} identifies the
specialization of the eigenvalue side, naturally in all its tensor
maps, and hence specializes the coherence diagrams themselves.
\end{proof}

\subsection{Betti comparison}
\label{found:betti}

Over $\mathbb C$ we use the realization of geometric adic constructible
complexes as algebraically constructible complexes of ordinary
$E$-vector spaces on the analytic space.  For a lisse sheaf it restricts
its continuous monodromy to topological loops.  On finite-type charts
it is conservative, preserves perversity, and commutes with the
pullbacks, tensor products, and finite-type pushes used below.  It also
identifies algebraic singular support with the complex conic analytic
singular support, as seen from comparison of vanishing-cycle tests.
These conventions and their passage to smooth stacks are those of
\cite[\S2.4]{CT}; see also \cite[\S2.2.1]{GR}.

We use Betti results to verify properties of already constructed adic
objects and adic comparison maps.  In particular, Betti perversity
bounds will justify adic truncation comparisons.  The infinite-rank
universal monodromic kernels used in the torus argument stay on the
Betti side.  Every object subsequently specialized is geometric adic
and lcc.
